% sec4-growth.tex -- Section 4 of the C3 paper (draft 2026-11-08): result D.
% Sources (run B, 2026-11-04): codex-20261104-x1-syzygies/03-FI-HOLE-AVOIDANCE.md and 04-NS-PC-PADDING.md
% (author x1, Codex); reviews notes-20261104-referee-x1-padding-c1.md (c1, Claude) and c2's review;
% adjudication codex-20261104-x2-padding-adjudication/01-ADJUDICATION.md (x2, Codex);
% cold referee notes-20261104/referee-runB-padding-transfer.md; Astra
% codex-20261104-runB-adjudication/00-ADJUDICATION.txt (C3) and 02-C3-SOURCE-AND-DEGREE.txt.
% Source excerpts: Rhoades Def 1.1 / Thm 3.12 (rhoades-source-excerpts.txt); Beame--Riis Def 8.2,
% Lemma 8.4, proof of Thm 8.1 (sources-20261101/beame-riis.txt).

\section{Growth in \texorpdfstring{$\ell$}{l}: the padding transfer}\label{sec:growth}

This section proves result D.  The mechanism is a transfer from Nullstellensatz designs to
PC lower bounds that costs $2D-2$ extra holes.  It is applied to the explicit designs of
Beame and Riis.

\begin{theorem}[padding transfer]\label{thm:padding}
Let $F$ be a field, $D\ge1$, $n_0\ge1$ and $k\ge0$.  If the board with $n_0+k$ pigeons and
$n_0$ holes has a unit degree-$D$ design over $F$, then for every
\[
  n\ \ge\ n_0+2D-2
\]
the board with $n+k$ pigeons and $n$ holes satisfies $\mathrm{ND}_D$ and $\mathrm{U}_D$.  In
particular $\mathrm{onto\text{-}FPHP}^{n+k}_n$ has no $\mathrm{PC}_F$ refutation of degree at
most $D$.
\end{theorem}

\begin{corollary}[D]\label{cor:growth}
Let $p$ be prime, $\ell\ge1$ and $F$ a field of characteristic $p$.  For
\[
  n\ \ge\ \frac{(p+2)^\ell-p^\ell}{2}+2^{\ell+1}-4,
\]
every $\mathrm{PC}_F$ refutation of $\mathrm{onto\text{-}FPHP}^{n+p^\ell}_n$ has degree at
least $2^\ell$.
\end{corollary}

\paragraph{Tier and inputs.}  T-WRITTEN.  The proof uses three published results that we do
not re-prove: Corollary~B of Church and Ellenberg~\cite{churchellenberg}, Theorem~3.12 of
Rhoades~\cite{rhoades2024}, and the design construction of Beame and
Riis~\cite[Def.~8.2, Lemma~8.4, proof of Thm.~8.1]{beame1998counting}.  It uses no
representation theory beyond these, and no result of Harman or of Hone et al.  The input we
consider most exposed is the conversion of Rhoades's theorem into filtered completeness on
the square (Lemma~\ref{lem:square-complete}); it is short, and we give it in full.
Section~\ref{sec:growth-provenance} records who wrote and checked each step.

\subsection{External inputs}\label{sec:growth-inputs}

\paragraph{FI-modules.}  An \emph{FI-module} assigns an abelian group $V(T)$ to each finite
set $T$ and a map $V(T)\to V(T')$ to each injection $T\hookrightarrow T'$, functorially.  It
is \emph{generated in degree $\le d$} if every $V(T)$ is spanned by the images of the
$V(T_0)$ with $|T_0|\le d$.  For a module $W$ that assigns a group to each finite set and
an action of bijections only, the \emph{induced} (``free'') FI-module is
\[
  M(W)(T)=\bigoplus_{T_0\subseteq T}W(T_0),
\]
with injections acting by relabelling~\cite[\S2.1, eq.~(2)]{churchellenberg}.  ``Free''
here does not mean projective, and no restriction on characteristic is involved.

\begin{theorem}[Church--Ellenberg~{\cite[Cor.~B]{churchellenberg}}]\label{thm:ce}
Let $M$ be an induced FI-module generated in degree $\le a$ and $V$ an FI-module generated
in degree $\le b$.  Then the kernel of every FI-module map $V\to M$ is generated in degree
$\le a+b+1$.
\end{theorem}

The theorem is stated for abelian groups.  We apply it to $F$-vector spaces by forgetting
the scalar structure; the images of smaller pieces are $F$-subspaces, so generation as
groups gives spanning over $F$.

\paragraph{The square.}  On an $s\times s$ board let $I_s\subseteq F[x]$ be generated by all
row sums $\sum_jx_{ij}$, all column sums $\sum_ix_{ij}$, and all products $x_{ij}x_{ij'}$,
$x_{ij}x_{i'j}$ of two variables in one row or one column, squares
included~\cite[Def.~1.1]{rhoades2024}.  Let $I(\mathcal P_s)$ be the ideal of polynomials
vanishing on the $s!$ permutation matrices, and $\operatorname{gr}I(\mathcal P_s)$ the ideal
of top-degree homogeneous parts of its elements.

\begin{theorem}[Rhoades~{\cite[Thm.~3.12]{rhoades2024}}]\label{thm:rhoades}
Over every field, $I_s=\operatorname{gr}I(\mathcal P_s)$.
\end{theorem}

Rhoades's later results on the module structure assume characteristic $0$ or $p>s$; we use
only Theorem~\ref{thm:rhoades}, which has no such assumption.

\paragraph{Beame--Riis designs.}  For $\ell\ge0$ put $N_\ell=((p+2)^\ell-p^\ell)/2$.
Beame and Riis construct, by induction on $\ell$, a $(2^\ell-1)$-design over
$\mathbb Z_p$ on the board with $N_\ell+p^\ell$ pigeons and $N_\ell$
holes~\cite[proof of Thm.~8.1]{beame1998counting}.  The step is their Lemma~8.4: a
$d$-design on a $D_0\times R_0$ board gives a $(2d+1)$-design on a board with
$(p+1)|D_0|+|R_0|$ pigeons and $|D_0|+(p+1)|R_0|$ holes.  Their $d$-designs
(Definition~8.2) are exactly our unit degree-$d$ designs with values in $\mathbb F_p$.  We
use the construction, not Theorem~8.1 itself; this avoids converting between their
multiplier degree and our product degree.

\subsection{Square filtered completeness}\label{sec:growth-square}

\begin{lemma}\label{lem:square-complete}
On an $s\times s$ board, for every $t$, $W_t$ is the set of elements of $F_t$ that vanish
at every permutation matrix.  Consequently every linear functional on $F_t$ that vanishes on
$W_t$ is the restriction of a linear combination, with coefficients in $F$, of evaluations at
permutation matrices.
\end{lemma}

\begin{proof}
Elements of $W_t$ vanish on permutation matrices, since each row or column relation does.
Conversely let $g\in F_t$ vanish on permutation matrices; lift it to a multilinear
polynomial $f$ of degree $\le t$ supported on matchings.  Then $f\in I(\mathcal P_s)$.  By
Theorem~\ref{thm:rhoades} its top homogeneous part, of degree $t'\le t$, is a combination
$\sum_r h_r\gamma_r$ of the generators $\gamma_r$ of $I_s$ with homogeneous $h_r$ and
$\deg(h_r\gamma_r)=t'$.  Let $\tilde\gamma_r$ be the corresponding axiom: $\gamma_r-1$
(a pigeon or hole axiom) for a row or column sum, $x^2-x$ (a Boolean axiom) for a square
$x^2$, and $\gamma_r$ itself (a collision or functional axiom) for a product of two distinct
variables in one line.  Each $\tilde\gamma_r$ has top part $\gamma_r$.  Then $f-\sum_rh_r\tilde\gamma_r$ has degree $<t'$ and still
vanishes on permutation matrices.  Repeating, $f$ is a combination of the axioms with
multipliers of degree at most $t-1$ for the linear axioms and $t-2$ for the quadratic ones.
Applying $\rho$, the Boolean, collision and functional terms vanish and each linear term becomes a combination of
row and column relations at bases of size $\le t-1$, so $g=\rho(f)\in W_t$.

For the second statement: a functional on $F_t$ vanishing on $W_t$ factors through the
evaluation map $F_t\to F^{\mathcal P_s}$, whose kernel is $W_t$.  Extend the induced
functional from the image to all of $F^{\mathcal P_s}$; every functional on
$F^{\mathcal P_s}$ is a combination of point evaluations.
\end{proof}

The functional obtained in this way need not be invariant, nonnegative or unital.  Every
evaluation at a permutation matrix $w$, viewed as the function $M\mapsto[M\subseteq w]$ on
matchings of every size, satisfies all row and column relations at every base: if
$M\subseteq w$ and $i$ is free in $M$, then $w(i)$ is free in $M$, so exactly one term of
the free row sum is $1$.  So do their linear combinations.

\subsection{Balanced arrays have small support}\label{sec:growth-support}

\begin{lemma}\label{lem:small-support}
Let $e\ge1$.  Every $f\in H_e(P,S)$ is a sum of arrays $f_\alpha\in H_e(P,S)$, each
supported on matchings inside $P_\alpha\times S_\alpha$ with $|P_\alpha|\le2e$ and
$|S_\alpha|\le2e$.
\end{lemma}

\begin{proof}
Fix the pigeon set $P$ and regard everything as a functor of the hole set.  Let $C_e(T)$ be
the space of arrays on $e$-matchings $P\leftrightarrow T$ all of whose free column sums
vanish.  Each column condition involves a single $e$-element hole set
$\mathrm{holes}(N)\cup\{j\}$, so
$C_e(T)=\bigoplus_{T_0\subseteq T,\,|T_0|=e}U(T_0)$, where $U(T_0)$ is the space of
column-balanced arrays using exactly the holes $T_0$.  With injections acting by
relabelling and extension by zero, $C_e$ is the induced module $M(U)$, generated in degree
$e$.  Let $B_{e-1}(T)$ have a basis of pairs $(N,i)$, $N$ an $(e-1)$-matching
$P\leftrightarrow T$ and $i$ a pigeon free in $N$; it is induced, generated in degree
$e-1$, since $(N,i)$ uses exactly the holes of $N$.  The row-sum map
$(Rf)(N,i)=\sum_{j}f(N\cup\{(i,j)\})$ is a map of FI-modules $C_e\to B_{e-1}$: after
extension by zero, the new terms vanish, and coordinates $(N,i)$ with $N$ outside the old
hole set get the value $0$ on both sides.  Its kernel is $H_e(P,\cdot)$.  By
Theorem~\ref{thm:ce} with $a=e-1$, $b=e$, the kernel is generated in degree
$\le 2e$: every $f\in H_e(P,S)$ is a sum of zero extensions of balanced arrays on hole sets
of size $\le2e$.  Extension by zero preserves both marginal conditions.

Apply the same argument with pigeons and holes exchanged to each summand, keeping its hole
set fixed.  This splits each summand further into balanced arrays on at most $2e$ pigeons,
without enlarging its hole set.
\end{proof}

\subsection{Extending one balanced layer}\label{sec:growth-tensor}

\begin{lemma}[square realisation]\label{lem:square-real}
Let $|A|=|B|=2e$ and $f\in H_e(A,B)$.  There is a function $\mu$ on the matchings of the
$A\times B$ board (of all sizes) that satisfies every row and column relation, vanishes on
matchings of size $<e$, and equals $f$ on matchings of size $e$.
\end{lemma}

\begin{proof}
The functional on $F_e$ that is $0$ below level $e$ and $f$ at level $e$ vanishes on $W_e$:
relations at bases of size $<e-1$ involve only zeros, and those at bases of size $e-1$ are
the vanishing marginals of $f$.  Lemma~\ref{lem:square-complete} and the remark after it
give $\mu$.
\end{proof}

\begin{lemma}[tensor extension]\label{lem:tensor}
Let $1\le e\le D-1$.  Split the board into pigeon sets $A\sqcup A'$ and hole sets
$B\sqcup B'$ with $|A|=|B|=2e$.  Let $\mu$ be as in Lemma~\ref{lem:square-real} on
$A\times B$ and let $\nu$ be a unit degree-$(D-e)$ design on $A'\times B'$.  For a matching
$M$ with $|M|\le D$ put
\[
  G(M)=\begin{cases}
  0 & \text{if $M$ has an edge in $A\times B'$ or $A'\times B$,}\\
  0 & \text{if $|M\cap(A\times B)|<e$,}\\
  \mu(M\cap(A\times B))\,\nu(M\cap(A'\times B')) & \text{otherwise.}
  \end{cases}
\]
Then $G$, as a functional on $F_D$, vanishes on $W_D$ and on all matchings of size $<e$,
and $G(M)=f(M)$ for matchings $M\subseteq A\times B$ of size $e$ (zero on other matchings of
size $e$).
\end{lemma}

\begin{proof}
$G$ is defined: in the third case the residual part has size at most $D-e$.  At level $e$
the residual part of a non-zero term is empty and $\nu(\emptyset)=1$.

Fix a base $M$ with $|M|\le D-1$ and a free pigeon $i$; columns are symmetric.  Write
$a=|M\cap(A\times B)|$ and $b=|M\cap(A'\times B')|$.  We check that the free row sum of $G$
at $(M,i)$ equals $G(M)$.
\begin{itemize}
\item $M$ has a crossing edge: every extension has one too; both sides are $0$.
\item $i\in A$, $a\le e-2$: extensions into $B$ have core size $\le e-1$ and extensions into
  $B'$ cross; both sides are $0$.
\item $i\in A$, $a=e-1$: $G(M)=0$.  Only extensions into $B$ can be non-zero, and their sum is
  $\nu(M\cap(A'\times B'))\cdot\sum_{j\in B\text{ free}}\mu(M_{\mathrm{core}}\cup\{(i,j)\})
  =\nu(\cdot)\,\mu(M_{\mathrm{core}})=0$, by the row relation of $\mu$ and $\mu=0$ below $e$.
  Here $b\le D-1-a=D-e$, so $\nu$ is evaluated within its degree.  This is the only place
  where the cutoff at level $e$ matters.
\item $i\in A$, $a\ge e$: the same row relation of $\mu$ gives the sum $G(M)$.
\item $i\in A'$, $a<e$: every extension into $B'$ still has core size $<e$, and extensions
  into $B$ cross; both sides are $0$.
\item $i\in A'$, $a\ge e$: then $b\le D-1-e$, so the row relation of $\nu$ at the base
  $M\cap(A'\times B')$ is one of its defining relations; extensions into $B$ cross; the sum
  is $\mu(M_{\mathrm{core}})\,\nu(M\cap(A'\times B'))=G(M)$.
\end{itemize}
So $G$ satisfies every relation at bases of size $\le D-1$, that is, $G(W_D)=0$.
\end{proof}

\begin{lemma}[padding a design]\label{lem:pad}
A unit degree-$D$ design on the board $(n_0+k)\times n_0$ gives one on $(r+k)\times r$ for
every $r\ge n_0$, and, by restriction, unit designs of every smaller degree.
\end{lemma}

\begin{proof}
Add $r-n_0$ pigeons and $r-n_0$ holes and fix a bijection $\sigma$ between them.  Define the
new design to be the old one times $[M_{\mathrm{new}}\subseteq\sigma]$ on matchings without
edges between old and new parts, and $0$ on others.  The proof of Lemma~\ref{lem:tensor}
applies with the old board as residual and $\sigma$ as core; there is no cutoff, and the
value at $\emptyset$ stays $1$.
\end{proof}

\subsection{Proof of Theorem~\ref{thm:padding}}\label{sec:growth-proof}

Let $n\ge n_0+2D-2$, and let $U$ be the unit degree-$D$ design on $(n+k)\times n$ given by
Lemma~\ref{lem:pad}.  Then $\mathrm{U}_D$ holds.

For $\mathrm{ND}_D$ we show that every functional $\lambda$ on $F_{D-1}$ vanishing on
$W_{D-1}$ extends to a functional on $F_D$ vanishing on $W_D$
(Section~\ref{sec:defs-written}).  Start with $\Lambda_0=\lambda(1)\,U$; the remainder
$\lambda-\Lambda_0|_{F_{D-1}}$ vanishes on $W_{D-1}$ and at $\emptyset$.  Suppose that for
some $1\le e\le D-1$ the remainder $r_e$, a functional on $F_{D-1}$ vanishing on $W_{D-1}$,
vanishes on all matchings of size $<e$.  Its level-$e$ layer $f_e$ lies in $H_e$ by the last
paragraph of Section~\ref{sec:defs-written} with $t=D-1$.  By
Lemma~\ref{lem:small-support}, $f_e=\sum_\alpha f_\alpha$ with each $f_\alpha$ balanced on at
most $2e$ pigeons and $2e$ holes.  Since $n\ge 2D-1>2e$, enlarge each support to exactly
$2e$ pigeons and $2e$ holes; the array stays balanced.  The complement has $n-2e\ge n_0$
holes and surplus $k$, so by Lemma~\ref{lem:pad} it has a unit degree-$(D-e)$ design.
Lemmas~\ref{lem:square-real} and~\ref{lem:tensor} give $G_\alpha$ vanishing on $W_D$ and
below level $e$, with level-$e$ layer $f_\alpha$.  Put
$\Lambda_e=\Lambda_{e-1}+\sum_\alpha G_\alpha$.  The new remainder
$r_{e+1}=r_e-\sum_\alpha G_\alpha|_{F_{D-1}}$ vanishes below level $e+1$ and still vanishes
on $W_{D-1}$.  After the step $e=D-1$ the remainder is zero, so $\Lambda_{D-1}$ extends
$\lambda$ and vanishes on $W_D$.  (For $D=1$ there are no correction steps.)

The PC conclusion is Proposition~\ref{prop:nd-pc}.  \qed

\subsection{Proof of Corollary~\ref{cor:growth}}\label{sec:growth-cor}

Take $D=2^\ell-1$, $n_0=N_\ell$ and $k=p^\ell$.  Beame and Riis's construction gives a unit
degree-$D$ design over $\mathbb F_p$ on $(n_0+k)\times n_0$; its defining conditions are
linear with coefficients in $\mathbb F_p$, so it is also a design over every field $F$ of
characteristic $p$.  For $\ell\ge1$, $n_0=N_\ell\ge N_1=1$.  Theorem~\ref{thm:padding}
excludes refutations of degree at most $2^\ell-1$ whenever
$n\ge N_\ell+2(2^\ell-1)-2=N_\ell+2^{\ell+1}-4$.  \qed

\begin{corollary}[exact degree at $p=2$, without E's inputs]\label{cor:growth-p2}
Let $\ell\ge1$.  Over every field of characteristic $2$, the PC degree of
$\mathrm{onto\text{-}FPHP}^{n+2^\ell}_n$ is exactly $2^\ell$ for
$n\ge(4^\ell-2^\ell)/2+2^{\ell+1}-4$.
\end{corollary}

\begin{proof}
At $p=2$ the bound $2^\ell$ of Corollary~\ref{cor:growth} equals the upper bound $p^\ell$ of
Proposition~\ref{prop:riis}.
\end{proof}

So at $p=2$ the exact degree of Corollary~\ref{cor:exact} has another proof, at a larger
threshold ($n\ge40$ instead of $15$ for $\ell=3$), using neither of E's
representation-theoretic inputs; both proofs use Riis's upper bound.  This corollary was
deduced while drafting the paper and checked in the model review of the paper, not in the
run that produced D.  For odd $p$ the two bounds differ, because Beame and Riis's designs stop at
degree $2^\ell-1$ while the Nullstellensatz degree may be as large as $p^\ell-1$; designs of
degree $p^\ell-1$, if they exist, would give the exact degree for every $p$ through this
section alone (Section~\ref{sec:open}).

For $\ell=1$ Corollary~\ref{cor:growth} gives degree at least $2$ for $k=p$ and every $n\ge1$, consistent with
Theorem~A.  For $p=2$, $\ell=2$ it gives degree at least $4$ for surplus $4$ from $n\ge10$;
the Lean result gives exactly $4$ from $n\ge5$.

\paragraph{The Nullstellensatz route, for comparison.}  Theorem~8.1 of
\cite{beame1998counting} states its bound in multiplier degree.  It also yields the
design used above: a certificate of product degree at most $D$, reduced modulo the Boolean
axioms, has every non-zero multiplier of degree at most $D-1$, since every axiom has positive
degree; with $D=2^\ell-1$ this contradicts Theorem~8.1, and duality gives the unit design.
The two routes give the same threshold; we cite the construction because it needs no
convention conversion.

\subsection{What was checked}\label{sec:growth-provenance}

The proof of Theorem~\ref{thm:padding} was written by a Codex agent (x1) in run 20261104-0759
(\texttt{03-FI-HOLE-AVOIDANCE.md}, \texttt{04-NS-PC-PADDING.md}).  Two Claude agents reviewed
it in the run (c1 re-did every case of Lemma~\ref{lem:tensor} by hand); a second Codex agent
adjudicated it against the primary sources of Church--Ellenberg and Rhoades; a separate
model referee with no contact with the run re-read both sources and found the transfer
consistent with its own small computations over $\mathbb F_2$ and $\mathbb F_3$.  The
adjudicator Astra checked the Beame--Riis construction in the author-hosted PDF
(Definitions~7.1, 7.2, 8.2, Lemma~8.4, proof of Theorem~8.1) and the conversion above.
The tensor construction of Lemma~\ref{lem:tensor} was also checked by exact computation at
$D=3$ over $\mathbb F_2$, with supports of size $2$ and $4$, on the full $12\times8$ board
(all $60{,}884$ free-line equations per case).  None of this is a Lean proof or a human
review, and none of it re-proves Corollary~B of \cite{churchellenberg}, Theorem~3.12 of
\cite{rhoades2024} or Lemma~8.4 of \cite{beame1998counting}.
